% Options for packages loaded elsewhere
\PassOptionsToPackage{unicode}{hyperref}
\PassOptionsToPackage{hyphens}{url}
\PassOptionsToPackage{dvipsnames,svgnames,x11names}{xcolor}
\documentclass[
  ngerman,
  10pt,
  a4paper,
]{article}
\usepackage{xcolor}
\usepackage[margin=22mm]{geometry}
\usepackage{amsmath,amssymb}
\setcounter{secnumdepth}{-\maxdimen} % remove section numbering
\usepackage{iftex}
\ifPDFTeX
  \usepackage[T1]{fontenc}
  \usepackage[utf8]{inputenc}
  \usepackage{textcomp} % provide euro and other symbols
\else % if luatex or xetex
  \usepackage{unicode-math} % this also loads fontspec
  \defaultfontfeatures{Scale=MatchLowercase}
  \defaultfontfeatures[\rmfamily]{Ligatures=TeX,Scale=1}
\fi
\usepackage{lmodern}
\ifPDFTeX\else
  % xetex/luatex font selection
  \setmainfont[Path=/Users/stefanhamann/.cache/codex-runtimes/codex-primary-runtime/dependencies/native/libreoffice-headless/libreoffice/LibreOfficeDev.app/Contents/Resources/fonts/truetype/,Extension=.ttf,BoldFont=DejaVuSans-Bold,ItalicFont=DejaVuSans-Oblique,BoldItalicFont=DejaVuSans-BoldOblique]{DejaVuSans}
  \setmonofont[Path=/Users/stefanhamann/.cache/codex-runtimes/codex-primary-runtime/dependencies/native/libreoffice-headless/libreoffice/LibreOfficeDev.app/Contents/Resources/fonts/truetype/,Extension=.ttf,BoldFont=DejaVuSansMono-Bold,ItalicFont=DejaVuSansMono-Oblique,BoldItalicFont=DejaVuSansMono-BoldOblique]{DejaVuSansMono}
\fi
% Use upquote if available, for straight quotes in verbatim environments
\IfFileExists{upquote.sty}{\usepackage{upquote}}{}
\IfFileExists{microtype.sty}{% use microtype if available
  \usepackage[]{microtype}
  \UseMicrotypeSet[protrusion]{basicmath} % disable protrusion for tt fonts
}{}
\makeatletter
\@ifundefined{KOMAClassName}{% if non-KOMA class
  \IfFileExists{parskip.sty}{%
    \usepackage{parskip}
  }{% else
    \setlength{\parindent}{0pt}
    \setlength{\parskip}{6pt plus 2pt minus 1pt}}
}{% if KOMA class
  \KOMAoptions{parskip=half}}
\makeatother
\ifLuaTeX
\usepackage[bidi=basic,shorthands=off]{babel}
\else
\usepackage[bidi=default,shorthands=off]{babel}
\fi
\ifPDFTeX
\else
\babelfont{rm}[Path=/Users/stefanhamann/.cache/codex-runtimes/codex-primary-runtime/dependencies/native/libreoffice-headless/libreoffice/LibreOfficeDev.app/Contents/Resources/fonts/truetype/,Extension=.ttf,BoldFont=DejaVuSans-Bold,ItalicFont=DejaVuSans-Oblique,BoldItalicFont=DejaVuSans-BoldOblique]{DejaVuSans}
\fi
\ifLuaTeX
  \usepackage{selnolig} % disable illegal ligatures
\fi
\setlength{\emergencystretch}{3em} % prevent overfull lines
\providecommand{\tightlist}{%
  \setlength{\itemsep}{0pt}\setlength{\parskip}{0pt}}
\usepackage{fvextra}
\RecustomVerbatimEnvironment{verbatim}{Verbatim}{breaklines=true,fontsize=\footnotesize}
\usepackage{xurl}
\usepackage{seqsplit}
\usepackage{adjustbox}
\usepackage{ragged2e}
\usepackage{titlesec}
\titleformat{\section}{\large\bfseries\raggedright}{}{0em}{}
\titleformat{\subsection}{\normalsize\bfseries\raggedright}{}{0em}{}
\DeclareRobustCommand{\codeword}[1]{\texttt{\seqsplit{#1}}}
\AtBeginEnvironment{longtable}{\small\RaggedRight}
\usepackage{newunicodechar}
\newunicodechar{𝒞}{\ensuremath{\mathcal C}}
\newunicodechar{𝔊}{\ensuremath{\mathfrak G}}
\newunicodechar{𝔤}{\ensuremath{\mathfrak g}}
\newunicodechar{𝓗}{\ensuremath{\mathcal H}}
\newunicodechar{ℋ}{\ensuremath{\mathcal H}}
\newunicodechar{𝔄}{\ensuremath{\mathfrak A}}
\newunicodechar{𝟙}{\ensuremath{\mathbf 1}}
\newunicodechar{𝕀}{\ensuremath{I}}
\newunicodechar{𝕋}{\ensuremath{\mathbb T}}
\newunicodechar{𝟘}{\ensuremath{\mathbf 0}}
\newunicodechar{⊞}{\ensuremath{\boxplus}}
\newunicodechar{∘}{\ensuremath{\circ}}
\newunicodechar{⟦}{\ensuremath{[\![}}
\newunicodechar{⟧}{\ensuremath{]\!]}}
\setlength{\emergencystretch}{4em}
\setlength{\parindent}{0pt}
\setlength{\parskip}{0.45em}
\AtBeginDocument{\hypersetup{colorlinks=true,linkcolor=black,urlcolor=blue}\pdfstringdefDisableCommands{\def\codeword#1{#1}}\RaggedRight}
\usepackage{fancyhdr}
\pagestyle{fancy}
\fancyhf{}
\fancyhead[L]{\small TFPT / Universalraum - v1.6.7}
\fancyfoot[R]{\thepage}
\setlength{\headheight}{15pt}
\usepackage{bookmark}
\IfFileExists{xurl.sty}{\usepackage{xurl}}{} % add URL line breaks if available
\urlstyle{same}
\hypersetup{
  pdflang={de-DE},
  colorlinks=true,
  linkcolor={Maroon},
  filecolor={Maroon},
  citecolor={Blue},
  urlcolor={Blue},
  pdfcreator={LaTeX via pandoc}}

\author{}
\date{}

\begin{document}

\section{TFPT / Universalraum: ein Vorgang, mehrere
Schatten}\label{tfpt-universalraum-ein-vorgang-mehrere-schatten}

\subsection{Einfach erklärt - v1.6.7, 15. September
2026}\label{einfach-erkluxe4rt---v1.6.7-15.-september-2026}

\textbf{Deine Schatten-Idee ist ein sinnvoller neuer Ansatz.} Ein
einziger Vorgang kann in verschiedenen Beschreibungen ganz
unterschiedlich aussehen. Wenn wir die Ansichten richtig verbinden,
können sie gemeinsam mehr verraten als jede für sich. Die heutige
Prüfung zeigt dafür einen konkreten Baustein im vorhandenen TFPT-Tensor.
Eine vollständige Erklärung unserer Welt haben wir damit noch nicht.

\subsection{Stell dir einen Gegenstand und mehrere Lampen
vor}\label{stell-dir-einen-gegenstand-und-mehrere-lampen-vor}

Eine Lampe wirft einen Kreis an die Wand. Daraus weißt du noch nicht, ob
der Gegenstand eine Kugel, eine Scheibe oder ein Zylinder ist. Eine
zweite Ansicht von der Seite kann diese Möglichkeiten unterscheiden.

Aber zwei Bedingungen sind entscheidend: Beide Lampen müssen wirklich
denselben Gegenstand beleuchten, und die zweite Ansicht muss etwas Neues
zeigen. Zwei Fotos desselben Schattens machen die Rekonstruktion nicht
vollständiger. Ebenso sind zwei Theorien nicht unabhängig bestätigt,
wenn dieselben Annahmen bereits in beide eingebaut wurden.

Für uns heißt das:

\begin{itemize}
\tightlist
\item
  \textbf{TFPT} liefert die innere Bau- und Symmetriestruktur.
\item
  \textbf{Universalraum} soll den gemeinsamen ausführbaren Prozess
  dahinter fassen.
\item
  \textbf{Die beobachtete Welt} liefert die Messungen, an denen dieser
  Prozess scheitern oder bestehen muss.
\end{itemize}

Das sind vorerst unterschiedliche Zugänge und Erkenntnisgrenzen. Wir
haben nicht nachgewiesen, dass sie drei physikalische Wände oder drei
getrennte Welten bilden. Holografie wäre eine stärkere Behauptung:
Information über ein Inneres müsste aus genau bestimmten
Randinformationen rekonstruierbar sein.

\subsection{Was daran jetzt wirklich konkret
ist}\label{was-daran-jetzt-wirklich-konkret-ist}

Im vorhandenen Tensor steckt schon eine Übersetzung: \textbf{Ein Zustand
mit 60 Komponenten lässt sich ohne Informationsverlust als Überlagerung
von Fermionpaaren schreiben.} Hier stehen zwei verschiedene
Darstellungen also tatsächlich für dieselbe Information.

Die Übersetzung sieht aber nicht jeden möglichen Paarzustand. Und wenn
eine einzelne Fermionmode verloren geht, kann bereits Information
fehlen. Ich habe das für alle 64 Moden nachgerechnet. Das ist deshalb
ein echter Kodierungsbaustein, aber noch kein geschütztes holografisches
Universum. Die Aufgabe lautet nun: Können die ursprünglichen Regeln
mehrere solcher Ansichten konsistent zusammenführen, einschließlich
ihrer Veränderungen?

\subsection{Drei einfache Verwechslungen haben den Blick
verstellt}\label{drei-einfache-verwechslungen-haben-den-blick-verstellt}

\textbf{Ein Name ist noch kein Ort.} Die 64 inneren Komponenten sagen
zunächst, welche Art von Zustand vorliegt. Sie sagen nicht, wo er liegt.
Auch mehrere Kopien derselben Art können verschiedene Zusammensetzungen
statt verschiedene Orte bedeuten. Raum muss aus unabhängig
adressierbaren Beziehungen entstehen.

\textbf{Gleichklang ist noch keine Übertragung.} Zwei passend
vorbereitete Uhren können zusammenhängende Anzeigen haben, ohne sich
gegenseitig anzutreiben. Das gilt auch für Quantenmodelle: Ich habe ein
exaktes Beispiel berechnet, in dem eine Kreuzkorrelation von null an
zeitabhängig wird, obwohl keine Signalübertragung möglich ist. Der
richtige Test ist: Ich greife nur bei A ein, und später ändert sich ohne
Aussortieren der Ergebnisse die Statistik bei B.

\textbf{Eine andere Zählweise ist noch kein neues Naturgesetz.} Derselbe
Tensor erlaubt zusätzliche einfache Paarerzeugungsprozesse. Manchmal
verschwinden diese aber wieder, wenn man andere zulässige
Teilchenvariablen benutzt. Wir können jetzt genau prüfen, wann das
innerhalb einer bestimmten kleinen Familie passiert. Das verhindert,
dass wir eine Umbenennung mit einer fundamentalen Entdeckung
verwechseln. Andere erlaubte Varianten sind wirklich verschiedene
Modelle; die Quelle muss zwischen ihnen auswählen.

\subsection{Wo wir mit der bisherigen Physik
stehen}\label{wo-wir-mit-der-bisherigen-physik-stehen}

Die Grundregel bleibt einfach: Zwei Fermionen werden in einen Vermittler
umgewandelt und zurück. Im festgelegten Modell ist ein eindeutiger
energetisch günstigster Zustand abgesichert. Entfernt man ein
ursprüngliches Fermion, trägt eine einzelne niedrige Energielinie mehr
als 88 Prozent des gesamten normierten Fermion-Spektralgewichts pro
Mode.

Bildlich: \textbf{Wir haben ein Instrument mit einem klaren Ton.} Noch
fehlt der Nachweis, wie dieser Ton auf derselben Grundlage zwischen
tatsächlich verschiedenen Teilen läuft und eine relativistische Welt
ergibt.

Diese Runde hat außerdem zwei begrenzte Strukturfragen gelöst: Die erste
Verzweigung besitzt genau vier symmetrische Richtungen, deren Kopplungen
nun berechnet sind. Und die 4035 bisher unsortierten bilinearen
Komponenten zerfallen in drei bestimmte Symmetrieklassen. Die richtige
kleine Rechenkette ist ebenfalls fünf Glieder weit exakt bekannt. Das
hilft bei der Rechnung, ist aber nicht der fehlende räumliche Ursprung.

\subsection{Was der übersehene Schlüssel sein
könnte}\label{was-der-uxfcbersehene-schluxfcssel-sein-kuxf6nnte}

Meine derzeit stärkste Arbeitshypothese lautet:

\textbf{Nicht noch ein größeres Objekt, sondern die gemeinsame Regel,
nach der Operationen und Beobachtungen dieselbe Ausführung beschreiben.}

Noethers Frage ist: Was bleibt bei erlaubten Veränderungen erhalten? Die
Schattenfrage ist: Welche Unterschiede können wir überhaupt erkennen?
Die Dynamikfrage ist: Welcher Eingriff kann wo etwas verändern? Diese
drei Fragen lassen sich gemeinsam am vorhandenen Tensor stellen. Das ist
eine Forschungsrichtung, keine bewiesene universelle Schlussformel.

\subsection{Die jetzt wichtigsten nächsten
Ziele}\label{die-jetzt-wichtigsten-nuxe4chsten-ziele}

\begin{enumerate}
\def\labelenumi{\arabic{enumi}.}
\tightlist
\item
  \textbf{Den richtigen einfachen Mechanismus auswählen.} Warum gilt
  genau die bisherige Paarregel und nicht eine andere ebenso
  symmetrische kleine Variante? Die nativen Operationen müssen diese
  Auswahl begründen.
\item
  \textbf{Eine gemeinsame Schattenkarte bauen.} Für denselben Zustand
  und dieselben Eingriffe ausrechnen, was jede Ansicht sieht und was sie
  übersieht. Eine Ansicht muss eine neue Vorhersage für die andere
  liefern, ohne dass wir hinterher Parameter passend machen.
\item
  \textbf{Einen echten kausalen Anschluss zeigen.} Ein verfügbarer
  Eingriff muss eine andere unabhängig festgelegte Messstelle
  beeinflussen. Erst darauf setzen Raumzeit, chirale Materie und
  Gravitation auf.
\end{enumerate}

Alle sechs neuen Texte sind geprüft und im Hauptdokument enthalten. Die
eigene Suite besteht aus 822 exakten Prüfbedingungen je normalem und
optimiertem Lauf; beide liefern dieselben Ergebnisdateien.
\textbf{Keines der acht großen TOE-Tore ist dadurch vollständig
geschlossen.} Der Fortschritt ist ein konkreter gemeinsamer
Kodierungsbaustein und ein wesentlich schärferer Test dafür, ob mehrere
Schatten wirklich denselben physikalischen Prozess rekonstruieren.

\end{document}
